\TieudeT{PHONG TỎA VẬT LÍ 11 - DAO ĐỘNG VÀ SÓNG}
\subsubsection*{PHẦN I. Thí sinh trả lời câu hỏi từ câu 1 đến câu 18. Mỗi câu hỏi thí sinh chỉ chọn một phương án}
\setcounter{ex}{0}
	\begin{ex}%[3P1Y1-1][Loai1]
		Một vật dao động điều hòa với biên độ A. Giá trị cực tiểu của li độ trong quá trình vật dao động là
		\choice
		{$A$}
		{$0$}
		{\True $-A$}
		{$-2A$}
		\loigiai{
			Giá trị cực tiểu của li độ trong quá trình vật dao động là $-A$.
		}
		%\haidong{6}
	\end{ex}
	\begin{ex}%[3P1Y1-2]
		Một vật nhỏ dao động điều hòa theo một quỹ đạo dài $24\ cm$. Dao động có biên độ
		\choice
		{$6\ cm$}
		{$3\ cm$}
		{\True $12\ cm$}
		{$24\ cm$}
		\loigiai{
			Quỹ đạo của dao động điều hòa bằng: $L = 2A =24\ cm \Rightarrow A = 12\  cm$.
		}	
		%\haidong{6}
	\end{ex}
	\begin{ex}%[3P1Y2-1][Loai1]
		Vận tốc của một chất điểm dao động điều hoà biến thiên
		\choice
		{cùng tần số và ngược pha với li độ}
		{\True cùng tần số và vuông pha với gia tốc}
		{khác tần số và vuông pha với li độ}
		{cùng tần số và cùng pha với li độ}
		\loigiai{
			Vận tốc của một chất điểm dao động điều hoà biến thiên cùng tần số và vuông pha với gia tốc.
		}
		%\haidong{5}
	\end{ex}
	%\loai{Thời điểm đầu tiên vật qua li độ cho trước tính cả chiều chuyển động}
	\begin{ex}%[3P1B8-2][Loai1]
		Một vật nhỏ dao động điều hòa có biên độ $8\ cm$, tần số góc $\dfrac{2\pi}{3}\ rad/s$, ở thời điểm ban đầu $t = 0$ vật qua vị trí có li độ $4\sqrt{3}\ cm$ theo chiều dương. Thời điểm đầu tiên kể từ $t = 0$ vật có li độ cực tiểu là
		\choice
		{$0,75\ s$}
		{$0,5\ s$}
		{\True $1,75\ s$}
		{$1,25\ s$}
		\loigiai{
			\begin{center}
				\begin{tikzpicture}
					\tkzDefPoints{0/0/o, 5/0/a, 4.325/0/b, 3.525/0/c, 2.5/0/d,  -5/0/h, -4.325/0/g, -3.525/0/f, -2.5/0/e, -6/0/x', 6/0/x}
					\tkzDefShiftPoint[b](-90:0.5cm){b1}
					\tkzDefShiftPoint[o](-90:1.2cm){o2}
					\tkzDefShiftPoint[a](-90:0.5cm){a1}
					\tkzDefShiftPoint[a](-90:1.2cm){a2}
					\tkzDefShiftPoint[h](-90:1.2cm){h2}
					\tkzDrawSegments[blue,->,very thick](x',x)
					\tkzDrawPoints[size=5,fill=black](o,a,b,c,d,e,f,g,h)
					\tkzLabelPoint[above](o){\tiny $O$}
					\tkzLabelPoint[above](a){\tiny$A$}
					\tkzLabelPoint[above](b){\tiny$\dfrac{A\sqrt{3}}{2}$}
					\tkzLabelPoint[above](c){\tiny$\dfrac{A\sqrt{2}}{2}$}
					\tkzLabelPoint[above](d){\tiny$\dfrac{A}{2}$}
					\tkzLabelPoint[above](h){\tiny$-A$}
					\tkzLabelPoint[above](g){\tiny $-\dfrac{A\sqrt{3}}{2}$}
					\tkzLabelPoint[above](f){\tiny$-\dfrac{A\sqrt{2}}{2}$}
					\tkzLabelPoint[above](e){\tiny$-\dfrac{A}{2}$}
					\tkzLabelPoint[above](x){\tiny$x$}
					\tkzDrawSegments[->,very thick,red](b1,a1 a2,h2)
					\tkzDrawSegments[very thick,red](a1,a2)
					\tkzLabelSegment[below](b1,a1){\tiny$\dfrac{T}{12}$}
					\tkzLabelSegment[below](a2,h2){\tiny$\dfrac{T}{2}$}
				\end{tikzpicture}
			\end{center}
			Vậy thời điểm cần tìm là: $t = \dfrac{T}{12}+\dfrac{T}{2}=\dfrac{7T}{12}=1,75\ s$.
		}
		%\haidong{10}
	\end{ex}
	\begin{ex}%[3P1K8-3][Loai1]
		Một chất điểm dao động điều hòa với chu kì $3\ s$, quỹ đạo là đoạn thẳng trên trục Ox quanh vị trí cân bằng là gốc tọa độ O. Biết rằng trong mỗi chu kì, khoảng thời gian để li độ của chất điểm lớn hơn $\sqrt{3}\ cm$ là $0,5\ s$. Tốc độ cực đại của vật trong quá trình dao động là
		\choice
		{$\dfrac{2\pi}{3}\ cm/s$}
		{$\dfrac{3\pi}{2}\ cm/s$}
		{\True $\dfrac{4\pi}{3}\ cm/s$}
		{$\dfrac{8}{3}\ cm/s$}
		\loigiai{
			\immini{\begin{itemize}
					\item 
					Ta có: $t = 0,5\ s = \dfrac{T}{6} \Rightarrow \Delta \varphi = \dfrac{\pi}{3}\ rad$.
					\item Ta biểu diễn các điểm có li độ $\sqrt{3}\ cm$ trên đường tròn lượng giác như hình vẽ.
					\item Chất điểm có li độ lớn hơn $\sqrt{3}\ cm$ khi nó đi từ vị trí $M_0$ qua biên dương, tới vị trí $M_1\Rightarrow A\dfrac{\sqrt{3}}{2} = \sqrt{3} \Rightarrow A = 2\ cm$ 
					$v_{\max} = 2.\dfrac{2\pi}{3} = \dfrac{4\pi}{3}\ cm/s$.
			\end{itemize}}{\begin{tikzpicture}[scale=1,thick,>=stealth']
					\tkzInit[ymin=-3,ymax=3,xmin=-3,xmax=3]\tkzClip
					\tkzDefPoints{-2.5/0/m, 2.5/0/n, 0/0/o, 0/2.5/p, 0/-2.5/q}
					\tkzDefShiftPoint[o](30:2){m'}
					\tkzDefShiftPoint[o](-30:2){0}
					\tkzDefPointBy[projection=onto m--n](m') \tkzGetPoint{h}
					\draw(o)circle(2cm);
					\tkzDrawSegments[dashed](m',0)
					\tkzDrawSegments(p,q)
					\tkzDrawSegments[->](m,n)
					\tkzDrawSegments[blue,->](o,0 o,m')
					\tkzDrawPoints[size=3](o,h,0,m')
					\tkzMarkAngles[size=0.7cm,arrows=->](0,o,m')
					\node at (0)[below right]{$M_0$};
					\node at (m')[above right]{$M_1$};
					\node at (2,0)[below right]{$A$};
					\node at (h)[below]{$\dfrac{\sqrt{3}}{2}A$};
			\end{tikzpicture}}
		}
		%\haidong{15}
	\end{ex}
	\begin{ex}%[3P1K9-2][Loai1]
		Một vật nhỏ dao động điều hòa theo phương trình $x=4\cos \left( 2t+\dfrac{\pi}{3} \right)$ cm (t tính bằng s). Khoảng thời gian ngắn nhất để vật đi từ vị trí có li độ $2\ cm$ đến vị trí có gia tốc $-8\sqrt{3}\ cm/s^2$ là
		\choice
		{$\dfrac{\pi }{6}\ s$}
		{$\dfrac{\pi}{24}\ s$}
		{$\dfrac{\pi}{8}\ s$}
		{\True $\dfrac{\pi}{12}\ s$}
		\loigiai{
			\begin{itemize}
				\item $a=-\omega^2x\Leftrightarrow x=2\sqrt{3}$ cm.
				\begin{center}
					\begin{tikzpicture}
						\tkzDefPoints{0/0/o, 5/0/a, 4.325/0/b, 3.525/0/c, 2.5/0/d,  -5/0/h, -4.325/0/g, -3.525/0/f, -2.5/0/e, -6/0/x', 6/0/x}
						\tkzDefShiftPoint[b](-90:0.5cm){b1}
						\tkzDefShiftPoint[d](-90:0.5cm){d1}
						\tkzDrawSegments[blue,->,very thick](x',x)
						\tkzDrawPoints[size=5,fill=black](o,a,b,c,d,e,f,g,h)
						\tkzLabelPoint[above](o){\tiny $O$}
						\tkzLabelPoint[above](a){\tiny$A$}
						\tkzLabelPoint[above](b){\tiny$\dfrac{A\sqrt{3}}{2}$}
						\tkzLabelPoint[above](c){\tiny$\dfrac{A\sqrt{2}}{2}$}
						\tkzLabelPoint[above](d){\tiny$\dfrac{A}{2}$}
						\tkzLabelPoint[above](h){\tiny$-A$}
						\tkzLabelPoint[above](g){\tiny $-\dfrac{A\sqrt{3}}{2}$}
						\tkzLabelPoint[above](f){\tiny$-\dfrac{A\sqrt{2}}{2}$}
						\tkzLabelPoint[above](e){\tiny$-\dfrac{A}{2}$}
						\tkzLabelPoint[above](x){\tiny$x$}
						\tkzDrawSegments[->,very thick,red](d1,b1)
					\end{tikzpicture}
				\end{center}
				\item Vậy thời gian ngắn nhất cần tìm là $\dfrac{T}{12}=\dfrac{\pi}{12}\ s$.
			\end{itemize}
		}
		%\haidong{15}
	\end{ex}
	\loai{Cho quãng đường viết phương trình dao động}
	
	%\loai{Quãng đường lớn nhất cần tách thời gian}
	\begin{ex}%[3P1KE-1][Loai1]
		Một vật nhỏ đang dao động điều hòa. Trong khoảng thời gian ngắn nhất giữa hai lần vật đi qua vị trí có li độ bằng $\dfrac{A}{2}$ liên tiếp vật đi được quãng đường dài nhất là
		\choice
		{$2A$}
		{$\dfrac{A}{2}$}
		{$\dfrac{A\sqrt{3}}{2}$}
		{\True $A\sqrt{3}$}
		\loigiai{
			\immini{
				\begin{itemize}
					\item 
					Ta biểu diễn trên đường tròn lượng giác như hình vẽ.
					\item Khoảng thời gian ngắn nhất giữa hai lần vật đi qua vị trí có li độ bằng $\dfrac{A}{2}$ liên tiếp ứng với khoảng thời gian vật đi từ $M_0$ qua biên dương tới $M_1$.\\
					$\Rightarrow \varphi = \dfrac{2\pi}{3} \Rightarrow \Delta t = \dfrac{T}{3} \Rightarrow \alpha = \dfrac{2\pi}{3}$.
					\item Ta có: $s_{\max} = 2A\sin\dfrac{\alpha}{2} = 2A\sin\dfrac{\pi}{3} = A\sqrt{3}$.
				\end{itemize}
			}{\begin{tikzpicture}[scale=1,thick,>=stealth']
					\tkzInit[ymin=-3,ymax=3,xmin=-3,xmax=3]\tkzClip
					\tkzDefPoints{-2.5/0/m, 2.5/0/n, 0/0/o, 0/2.5/p, 0/-2.5/q}
					\tkzDefShiftPoint[o](60:2){0}
					\tkzDefShiftPoint[o](-60:2){m'}
					\tkzDefPointBy[projection=onto m--n](m') \tkzGetPoint{h}
					\draw(o)circle(2cm);
					\tkzDrawSegments[dashed](m',0)
					\tkzDrawSegments(p,q)
					\tkzDrawSegments[->](m,n)
					\tkzDrawSegments[blue,->](o,0 o,m')
					\tkzDrawPoints[size=3](o,h,0,m')
					\tkzMarkAngles[size=0.7cm,arrows=->](m',o,0)
					\node at (m')[below right]{$M_0$};
					\node at (0)[above right]{$M_1$};
					\node at (2,0)[below right]{$A$};
					\node at (h)[below right]{$\dfrac{A}{2}$};
			\end{tikzpicture}}
		}
		%\haidong{15}
	\end{ex}
	
	
	\begin{ex}%[3P1-19.1K][Loai1] 
		\immini[thm]{
			Một vật dao động điều hòa có đồ thị như hình vẽ. Phương trình dao động của vật là
			\choice
			{$x=8\cos \left( \dfrac{\pi}{5}t-\dfrac{5\pi}{6} \right)\ cm$}
			{$x=8\cos \left( \dfrac{\pi}{5}t+\dfrac{5\pi}{6} \right)\ cm$}
			{$x=8\cos \left( \dfrac{3\pi}{10}t-\dfrac{3\pi}{4} \right)\ cm$}
			{\True $x=8\cos \left( \dfrac{3\pi}{10}t+\dfrac{3\pi}{4} \right)\ cm$}}{
			\begin{tikzpicture}[draw=Mapcolor,scale=1,thick,>=stealth']
				\pgfplotsset{width=8cm,height=4cm,compat=1.4}
				\begin{axis}[
					axis lines=middle,% Tùy chọn hệ trục tọa độ Đề các
					xmin=0,xmax=5.5,xstep=1,ymin=-2.16,ymax=2.5,ystep=1,
					grid=major, %Có các tùy chọn: minor|major|both|none
					x grid style={dashed,Mapcolor},
					y grid style={dashed,Mapcolor},
					ytick={-2,2},
					xtick={0.5,1,...,5},
					xticklabels={$\frac{5}{6}$,,,,$\frac{25}{6}$},
					yticklabels={$-8$,$+8$},
					xticklabel style={Mapcolor,above=3pt},
					xlabel style={right=3pt},
					xlabel=$t(s)$,
					ylabel style={above=3pt},
					ylabel={$x\ (cm)$}
					]
					\addplot[line width=1.2pt,domain=0:5,samples=200,Mapcolor]{2*cos(deg(0.5*pi*x+(3/4)*pi))};
				\end{axis}
			\end{tikzpicture}
		}
		\loigiai{
			\begin{itemize}
				\item Từ đồ thị, ta có: $T=2\left( \dfrac{25}{6}-\dfrac{5}{6} \right)=\dfrac{20}{3}\Rightarrow \omega =\dfrac{3\pi}{10}$ rad/s.
				\item Biên độ dao động của vật $A = 8$ cm.
				\item Thời điểm $t=\dfrac{5}{6}$ s, vật đi qua vị trí biên âm, thời điểm $t = 0$ tương ứng với góc lùi $\Delta\varphi=\omega\Delta t=0,25\pi\ rad$.
				\item Vậy $x=8\cos \left( \dfrac{3\pi}{10}t+\dfrac{3\pi}{4} \right)$ cm.
			\end{itemize}
		}
		%\haidong{10}
	\end{ex}
	\begin{ex}%[3P5-3.2K][Loai9]
		Trong thí nghiệm Young về giao thoa ánh sáng, hai khe được chiếu bằng ánh sáng đơn sắc có bước sóng $0,6\mu m$. Khoảng cách giữa hai khe sáng là $1\ mm$, khoảng cách từ mặt phẳng chứa hai khe đến màn quan sát là $1,5\ m$. Trên màn quan sát, khoảng cách từ vân sáng bậc $2$ đến vân tối thứ $10$ khác phía so với vân trung tâm cách nhau một đoạn là
		\choice
		{$6,45\ mm$}
		{\True $10,35\ mm$}
		{$10,9\ mm$}
		{$8,8\ mm$}
		\loigiai{
			\begin{itemize}
				\item Khoảng vân giao thoa trên màn là \bth{i=\dfrac{\lambda.D}{a}=\dfrac{0,6.1,5}{1}=0,9\ mm}.
				\begin{itemize}
					\item Vị trí vân sáng bậc 2: \bth{x_{s(2)}=2i=2.0,9=1,8\ mm}
					\item Vị trí vân tối thứ 10: \bth{x_{t(10)}=\pbrace{9+\dfrac{1}{2}}i=\pbrace{9+\dfrac{1}{2}}.0,9=8,55\ mm}
				\end{itemize}
				\item Vì hai vân khác phía nên khoảng cách giữa hai vân là
				\begin{align*}
					\Delta x=x_{s(2)}+x_{t(10)}=1,8+8,55=10,35\ mm.
				\end{align*}
			\end{itemize}
		}
	\end{ex}
	\begin{ex}%[3P2BB-1][Loai1]
		Một âm có tần số xác định truyền lần lượt trong nhôm, nước, không khí với tốc độ tương ứng là $v_1,\ v_2,\ v_3$. Nhận định nào sau đây là đúng?
		\choice
		{$v_2 > v_1 > v_3$}
		{$v_3 > v_2 > v_1$}
		{$v_1 > v_3 > v_2$}
		{\True $v_1 > v_2 > v_3$}
		%	\loigiai{
			%		Tốc độ truyền âm giảm dần từ chất rắn đến lỏng đến khí.
			%	}
	\end{ex}
	\begin{ex}%[3P2BB-1][Loai1]
		Khi nghe hai ca sĩ hát ở cùng một độ cao, ta vẫn phân biệt được giọng hát của từng người là do
		\choice
		{tần số và cường độ âm khác nhau}
		{\True âm sắc của mỗi người khác nhau}
		{tần số và năng lượng âm khác nhau}
		{tần số và biên độ âm khác nhau}
		%	\loigiai{
			%		Ta phân biệt được hai âm ở cùng một độ cao (tần số) là do âm sắc của mỗi âm là khác nhau.
			%	}
	\end{ex}
	\begin{ex}
		Khi nói về sóng điện từ, phát biểu nào sau đây là \textbf{sai}?
		\choice
		{Sóng điện từ truyền trong chân không với tốc độ $3.10^8\ m/s$}
		{Sóng điện từ là sóng ngang}
		{\True Sóng điện từ chỉ truyền được trong môi trường rắn, lỏng, khí}
		{Sóng điện từ có thể bị phản xạ khi gặp mặt phân cách giữa hai môi trường}
		%	\loigiai{
			%		Sóng điện từ truyền được trong rắn, lỏng, khí và cả trong chân không}
	\end{ex}
	\begin{ex}%[3P2Y1-1][Loai1]
		Phát biểu nào sau đây về đại lượng đặc trưng cho sóng cơ học là \textbf{không} đúng?
		\choice
		{Chu kì của sóng đúng bằng chu kì dao động của các phần tử môi trường}
		{Bước sóng là quãng đường sóng truyền đi được trong một chu kì}
		{\True Tốc độ truyền sóng đúng bằng tốc độ dao động của các phần tử môi trường}
		{Tần số của sóng đúng bằng tần số dao động của các phần tử môi trường}
		%	\loigiai{Trong sóng cơ, tốc độ truyền sóng là tốc độ truyền pha dao động, không phải là tốc độ dao động của các phần tử sóng}
	\end{ex}
	\begin{ex}%[3P5Y8-1][Loai1]
		Nguồn bức xạ nào sau đây \textbf{không} phát ra tia tử ngoại?	
		\choice
		{Mặt trời}
		{Đèn hơi thủy ngân}
		{Hồ quang điện}
		{\True Ngọn nến}	
		%	\loigiai{
			%		Ngọn nến không phát ra tia tử ngoại.
			%	}
	\end{ex}
	\begin{ex}%[3P2B5-2][Loai1] 
		Trong thí nghiệm tạo vân giao thoa sóng trên mặt nước, người ta dùng nguồn dao động có tần số $50\ Hz$ và đo được khoảng cách giữa hai cực đại liên tiếp nằm trên đường nối hai tâm dao động là $2\ mm$. Tốc độ truyền sóng trên dây là
		\choice
		{$40\ cm/s$}
		{$10\ cm/s$}
		{\True $20\ cm/s$}
		{$30\ cm/s$}
		\loigiai{
			\begin{itemize}
				\item Khoảng cách giữa hai cực đại liên tiếp là $\dfrac{\lambda}{2} = 2\ mm\ \Rightarrow \lambda =4\ mm \Rightarrow v = \lambda .f = 50.4 = 200 = 20$ cm/s.
				\item 
			\end{itemize}
		}
	\end{ex}
	\begin{ex}%[3P5-3.2K][Loai8]
		Tiến hành thí nghiệm Young về giao thoa ánh sáng với ánh sáng đơn sắc có bước sóng $0,6\ \mu m$. Khoảng cách giữa hai khe là $0,3\ mm$, khoảng cách từ mặt phẳng chứa hai khe đến màn quan sát là $2\ m$. Trên màn, khoảng cách giữa vân sáng bậc $3$ và vân sáng bậc $5$ ở hai phía so với vân sáng trung tâm là
		\choice
		{$8\ mm$}
		{\True $32\ mm$}
		{$20\ mm$}
		{$12\ mm$}
		\loigiai{
			\begin{itemize}
				\item Do hai vân sáng khác phía: $\Delta x=\dfrac{\lambda D}{a}\left(k_5+k_3\right)=\dfrac{0,6.10^{-6}.2}{0,3.10^{-3}}(5+3)=0,032\ m=32\ mm$.
				\item 
				\item 
			\end{itemize}
		}
	\end{ex}
	\begin{ex}%[3P2G6-2][Loai1]
		Trên mặt nước, hai nguồn kết hợp A, B cách nhau $20\ cm$ dao động điều hòa cùng pha, tạo ra sóng có bước sóng $3\ cm$. Xét các điểm trên mặt nước thuộc đường tròn tâm A, bán kính AB, điểm nằm trên đường tròn dao động với biên độ cực đại cách xa đường trung trực của AB nhất một khoảng bằng
		\choice
		{$34,5\ cm$}
		{\True $26,1\ cm$}
		{$21,7\ cm$}
		{$19,7\ cm$}	
		\loigiai{
			\immini{\begin{itemize}
					\item $\text{số nguyên lớn nhất}<\dfrac{AB}{\lambda}=\dfrac{20}{3}=6,7\Rightarrow n=6$.
					\item Điểm M phải là cực đại gần A nhất nên:\\ $MA-MB=-6\lambda \Rightarrow MB=38\ cm$.
					\item $\cos \alpha =\dfrac{AB^2+MB^2-MA^2}{2.AB.MB}=0,95$.
			\end{itemize}}{\begin{tikzpicture}[scale=0.7,thick,>=stealth']
					\tkzInit[ymin=-4,ymax=4,xmin=-3.5,xmax=4]\tkzClip
					\tkzDefPoints{0/0/a, 1.5/0/o, 3/0/b, 1.5/3.5/p, 1.5/-3.5/q, -0.5/3/y, -0.5/-3/y', -3/0/b'}
					\tkzDefShiftPoint[a](100:3){m}
					\tkzDefShiftPoint[a](-100:3){n}
					\tkzDefPointBy[projection=onto p--q](m) \tkzGetPoint{i}
					\tkzDefPointBy[projection=onto b--b'](m) \tkzGetPoint{h}
					\tkzDrawCircle[radius](a,b)
					\tkzDrawSegments(b',b b,m)
					\tkzDrawSegments[dashed](h,m i,m p,q)
					\tkzDrawPoints[size=3](a,b,o,m,n)
					\tkzMarkAngles[size=0.5cm](m,b,b')
					\tkzLabelAngle[pos=0.8](m,b,b'){$\alpha$}
					\draw[blue](y) to [bend left =30](y');
					\node at (a)[below]{$A$};
					\node at (h)[below]{$H$};
					\node at (i)[right]{$I$};
					\node at (b)[below right]{$B$};
					\node at (m)[above right]{$M$};
					\node at (o)[below left]{$O$};
		\end{tikzpicture}}}
	\end{ex}
	
	
	\begin{ex}%[3P2K8-1][Loai1]
		Trong một thí nghiệm về sóng dừng trên một sợi dây có chiều dài $\ell$ với hai đầu cố định. Trên dây có sóng dừng với ba bụng sóng. Biết tốc độ truyền sóng trên dây là $v$. Thời gian giữa hai lần sợi dây duỗi thẳng liên tiếp là
		\choice
		{\True $\dfrac{\ell}{3v}$}
		{$\dfrac{2\ell}{3v}$}
		{$\dfrac{v}{3\ell}$}
		{$\dfrac{2v}{3\ell}$}
		%	\loigiai{
			%		\begin{itemize}
				%			\item $\lambda = \dfrac{2\ell}{3}\Rightarrow$ $T = \dfrac{\lambda}{v} = \dfrac{2\ell}{3v}$.
				%			\item Thời gian giữa 2 lần sợi dây duỗi thẳng liên tiếp là:
				%			$\dfrac{T}{2} = \dfrac{\ell}{3v}$.
				%		\end{itemize}
			%	}
	\end{ex}

\subsubsection*{PHẦN 2. Thí sinh trả lời từ câu 1 đến câu 4. Trong mỗi ý a), b), c), d) ở mỗi câu, thí sinh chọn đúng hoặc sai.}
\setcounter{ex}{0}
\begin{ex}
	Trong dao động điều hòa, động năng và thế năng của vật biến thiên theo thời gian và có mối liên hệ chặt chẽ với chuyển động của vật.
	\choiceTF[t]
	{\True Động năng đạt giá trị cực đại khi gia tốc đổi chiều}
	{\True Thế năng đạt giá trị cực đại khi vật đổi chiều chuyển động}
	{Động năng đạt giá trị cực tiểu khi vật đi qua vị trí cân bằng theo chiều âm của trục $Ox$}
	{\True Thế năng đạt cực tiểu khi vật chuyển từ chuyển động nhanh dần sang chậm dần}
	\loigiai{
		\begin{itemchoice}
			\itemch Đúng.
			\itemch Đúng.
			\itemch Sai.
			\itemch Đúng.
		\end{itemchoice}
	}
	%\haidong{25}
\end{ex}


\begin{ex}
	Cho một chất điểm dao động điều hòa với biên độ bằng $4 \ cm$, chu kì bằng $0,5 \ s$.
	\choiceTF[t]
	{Trong mỗi chu kì, khoảng thời gian mà tốc độ chuyển động của vật lớn hơn $8 \sqrt{2} \pi\ cm/s$ là $0,75 \ s$}
	{Trong mỗi chu kì, khoảng thời gian mà li độ của vật lớn hơn $3 \ cm$ là $0,4 \ s$}
	{\True Ở thời điểm $t$ vật chuyển động theo chiều dương qua li độ $4 \ cm$. Sau lúc đó $\dfrac{1}{24}\ s$, li độ và chiều chuyển động của vật là $x=2 \sqrt{3}\ cm$ và chuyển động theo chiều âm}
	{\True Khoảng thời gian ngắn nhất giữa hai lần liên tiếp vật đạt tốc độ $8 \pi\ cm/s$ là $\dfrac{1}{12}\ s$}
	\loigiai{
		\begin{itemchoice}
			\itemch Sai.
			\itemch Sai.
			\itemch Đúng.
			\itemch Đúng.
		\end{itemchoice}
	}
	%\haidong{26}
\end{ex}
\begin{ex}
	Dựa vào phương dao động của các phần tử và phương lan truyền của sóng người ta phân sóng thành hai loại là sóng dọc và sóng ngang.
	\choiceTF[t]
	{\True Sóng dọc là sóng trong đó có phương dao động (của các phần tử môi trường) trùng với phương truyền}
	{Sóng ngang là sóng trong đó có phương dao động (của các phần tử môi trường) trùng với phương truyền} 
	{\True Sóng dọc truyền được trong các môi trường chất rắn và bề mặt chất lỏng}
	{Sóng ngang truyền được trong môi trường rắn, lỏng, khí}
	\loigiai{
		\begin{itemchoice}
			\itemch Đúng.
			\itemch Sai. Sóng ngang là sóng trong đó có phương dao động (của các phần tử môi trường) vuông góc với phương truyền
			\itemch Đúng. Sóng dọc truyền được trong chất rắn, lỏng, khí
			\itemch Sai. Sóng ngang truyền được trong chất rắn và trên bề mặt chất lỏng
		\end{itemchoice}
	}
\end{ex}

\begin{ex}
	Trong thí nghiệm giao thoa Young về giao thoa ánh sáng có bước sóng $\lambda$, khoảng cách từ màn tới hai khe là $1,5 \ m$, khoảng cách giữa hai khe là $a=1 \ mm$. Biết khoảng cách từ vân trung tâm tới vân sáng bậc $5$ là $4,5 \ mm$.
	\choiceTF[t]
	{\True Khoảng vân là $i=0,9\ mm$}
	{Bước sóng giao thoa là $\lambda=6000\ nm$}
	{Điểm cách vân trung tâm $3,5 \ mm$ là vân sáng bậc $4$}
	{Giá sử thay đổi nguồn sáng phát ra bước sóng $\lambda=0,4\ \mu m$ thì khoảng vân tăng $1,5$ lần}
	\loigiai{
		\begin{itemchoice}
			\itemch Đúng.
			\itemch Sai.
			\itemch Sai.
			\itemch Sai.
		\end{itemchoice}
	}
\end{ex}
\subsubsection*{PHẦN 3. Thí sinh trả lời từ câu 1 đến câu 6.}
\setcounter{ex}{0}
\begin{ex}
	Một chất điểm dao động điều hoà với phương trình li độ $x=2 \cos \left(\pi t+\dfrac{\pi}{6}\right)\ cm$ ($t$ tính bằng giây). Lấy $\pi^2=10$. Tốc độ của chất điểm ở thời điểm $t=0,5 \ s$ bằng bao nhiêu $cm/s$ (làm tròn kết quả đến chữ số hàng phần trăm)?
	\shortans[oly]{5,44}
	\loigiai{
		$5,44 \ cm/s$
	}
	%\haidong{30}
\end{ex}
\begin{ex}
	Để đo chiều dài của một dãy phòng học, do không có thước để đo trực tiếp, nên một học sinh đã làm như sau: Lấy một cuộn dây chỉ mảnh, không giãn, căng và đo lấy một đoạn bằng chiều dài của dãy phòng, sau đó gấp đoạn chỉ đó làm 74 phần bằng nhau. Dùng một con lắc đơn có chiều dài dây treo bằng chiều dài của một phần vừa gấp, kích thích cho con lắc dao động với biên độ góc nhỏ thì thấy con lắc thực hiện được 10 dao động toàn phần trong 18 giây. Lấy $g=9,8\ m/s^2$. Dãy phòng học mà bạn học sinh đo được là bao nhiêu m (làm tròn kết quả đến chữ số hàng đơn vị)? 
	\shortans[oly]{60}
	\loigiai{
		\begin{itemize}
			\item $\Delta t=10T=18\ s\Rightarrow T=1,8\ s=2\pi\sqrt{\dfrac{\ell}{g}}\Rightarrow \ell\approx 0,8043\ m\Rightarrow L=74.0,8043\approx 60\ m$
		\end{itemize}
	}	
	\haidong{15}
	
	
\end{ex}
\begin{ex}
	Một vật nhỏ dao động điều hòa trên trục $Ox$ với biên độ bằng $4 \ cm$ và tần số bằng $2 \ Hz$. Khoảng thời gian ngắn nhất tính từ thời điểm vật có vận tốc bằng $8 \pi\ cm/s$ đến thời điểm vật có gia tốc bằng $32 \sqrt{2} \pi^2\ cm/s^2$ là bao nhiêu giây (làm tròn kết quả đến chữ số hàng phần trăm)?
	\shortans[oly]{0,23}
	\loigiai{
		0,23
	}
	%\haidong{30}
\end{ex}
\begin{ex}
	Trong một thí nghiệm tạo vân giao thoa trên sóng nước, người ta dùng hai nguồn dao động đồng pha có bước sóng $10 \ cm$. Khoảng cách giữa $6$ điểm cực tiểu liên tiếp trên đoạn nối hai nguồn là bao nhiêu cm?
	\shortans[oly]{25}
	\loigiai{
		$25 \ cm$
	}
\end{ex}
\begin{ex}
	Tiến hành thí nghiệm Young về giao thoa ánh sáng với ánh sáng đơn sắc nguồn sáng phát ra ánh sáng đơn sắc có bước sóng $\lambda$. Trên màn quan sát, vân sáng bậc $5$ cách vân sáng trung tâm $4,5\ mm$. Khoảng vân giao thoa trên màn bằng bao nhiêu milimét?
	\shortans[oly]{0,9}
	\loigiai{
		Vân sáng bậc $5: x=5 i \Rightarrow i=\dfrac{x}{5}=0,9 \ mm$.
		Đáp số: 0,9
	}
\end{ex}
\begin{ex}
	Trong thí nghiệm Young về giao thoa ánh sáng, khoảng cách giữa hai khe là $2 \ mm$, khoảng cách từ mặt phẳng chứa hai khe đến màn quan sát là $1,2 \ m$. Trên màn, quan sát được $9$ vân sáng mà khoảng cách giữa hai vân sáng ngoài cùng là $2,4 \ mm$. Bước sóng của ánh sáng trong thí nghiệm bằng bao nhiêu $\mu m$?
	\shortans[oly]{0,5}
	\loigiai{
		\begin{itemize}
			\item Khoảng cách 9 vân sáng liên tiếp: $8 i=2,4 \ mm \Rightarrow i=0,3 \ mm$
			\item 		Bước sóng: $\lambda=\dfrac{a \cdot i}{D}=\dfrac{2 \cdot 0,3}{1,2}=0,5\ \mu m$
		\end{itemize}
		Đáp số: 0,5
	}
\end{ex}